% PTIT-style preliminary pages. Do not reuse identities or statements from sample group.
\chapter*{LỜI CAM ĐOAN}
\addcontentsline{toc}{chapter}{Lời cam đoan}

Chúng em, nhóm sinh viên thực hiện, gồm có:

\begin{enumerate}
\setlength{\itemsep}{4pt}
\setlength{\parsep}{0pt}
\item Dương Ngọc Phú
\item Phạm Thế Vỹ
\item Vũ Văn Kiên
\item Lê Hoàng Nguyên Quân
\end{enumerate}

Chúng em xin cam đoan đề tài “Nghiên cứu Ensemble Learning cho bài toán phát hiện xâm nhập mạng: Thực nghiệm phân loại nhị phân trên UNSW-NB15” là công trình nghiên cứu do tập thể nhóm thực hiện dưới sự hướng dẫn của TS. Lê Ngọc Hiếu.

Các số liệu, quy trình và kết quả thực nghiệm trình bày trong báo cáo được nhóm ghi nhận trung thực, phù hợp với dữ liệu và chương trình thực nghiệm đã sử dụng. Những nội dung tham khảo từ các nguồn khác đều được trích dẫn và ghi rõ nguồn gốc.

Nhóm sinh viên chịu trách nhiệm về tính chính xác và tính nguyên gốc của nội dung báo cáo, đồng thời công khai các giới hạn về dữ liệu, thiết kế kiểm định và khả năng tái lập để người đọc đánh giá đúng phạm vi kết luận.

\vspace{1.8cm}
\begin{flushright}
TP. Hồ Chí Minh, năm 2026\\[10pt]
Nhóm sinh viên thực hiện
\end{flushright}

\clearpage
\chapter*{LỜI CẢM ƠN}
\addcontentsline{toc}{chapter}{Lời cảm ơn}
Nhóm sinh viên trân trọng cảm ơn giảng viên hướng dẫn TS. Lê Ngọc Hiếu đã đồng hành trong quá trình thực hiện đề tài nghiên cứu khoa học về Ensemble Learning cho phát hiện xâm nhập mạng.

Chúng em cảm ơn các thầy cô của Khoa Công nghệ Thông tin 2, Học viện Công nghệ Bưu chính Viễn thông, vì nền tảng kiến thức và môi trường học thuật hỗ trợ quá trình học tập, nghiên cứu. Báo cáo được thực hiện dựa trên những dữ liệu và nguồn tham khảo có ghi nhận, với mong muốn làm rõ chất lượng thực nghiệm cũng như những hạn chế khi ứng dụng các mô hình học máy cho NIDS.

Nhóm mong nhận được những ý kiến nhận xét và phản biện để tiếp tục hoàn thiện nghiên cứu.

\vspace{1.8cm}
\begin{flushright}
TP. Hồ Chí Minh, năm 2026\\[10pt]
Nhóm sinh viên thực hiện
\end{flushright}
